\documentclass[10pt,a4paper]{amsart}
\usepackage{amsmath,amssymb,amsthm,graphicx}
\usepackage{verbatim}
\newenvironment{reference}{\comment}{\endcomment}
%\newenvironment{reference}{}{}
\newenvironment{slogan}{\comment}{\endcomment}
\newenvironment{history}{\comment}{\endcomment}

% For commutative diagrams we use Xy-pic
\usepackage[all]{xy}

% We use 2cell for 2-commutative diagrams.
\xyoption{2cell}
% Theorem environments.
%
\theoremstyle{plain}
\newtheorem{theorem}[subsection]{Theorem}
\newtheorem{proposition}[subsection]{Proposition}
\newtheorem{lemma}[subsection]{Lemma}

\theoremstyle{definition}
\newtheorem{definition}[subsection]{Definition}
\newtheorem{example}[subsection]{Example}
\newtheorem{exercise}[subsection]{Exercise}
\newtheorem{situation}[subsection]{Situation}

\theoremstyle{remark}
\newtheorem{remark}[subsection]{Remark}
\newtheorem{remarks}[subsection]{Remarks}

\numberwithin{equation}{subsection}

% Macros
%
\def\lim{\mathop{\mathrm{lim}}\nolimits}
\def\colim{\mathop{\mathrm{colim}}\nolimits}
\def\Spec{\mathop{\mathrm{Spec}}}
\def\Hom{\mathop{\mathrm{Hom}}\nolimits}
\def\Ext{\mathop{\mathrm{Ext}}\nolimits}
\def\SheafHom{\mathop{\mathcal{H}\!\mathit{om}}\nolimits}
\def\SheafExt{\mathop{\mathcal{E}\!\mathit{xt}}\nolimits}
\def\Sch{\mathit{Sch}}
\def\Mor{\mathop{\mathrm{Mor}}\nolimits}
\def\Ob{\mathop{\mathrm{Ob}}\nolimits}
\def\Sh{\mathop{\mathit{Sh}}\nolimits}
\def\NL{\mathop{N\!L}\nolimits}
\def\CH{\mathop{\mathrm{CH}}\nolimits}
\def\proetale{{pro\text{-}\acute{e}tale}}
\def\etale{{\acute{e}tale}}
\def\QCoh{\mathit{QCoh}}
\def\Ker{\mathop{\mathrm{Ker}}}
\def\Im{\mathop{\mathrm{Im}}}
\def\Coker{\mathop{\mathrm{Coker}}}
\def\Coim{\mathop{\mathrm{Coim}}}

% Boxtimes
%
\DeclareMathSymbol{\boxtimes}{\mathbin}{AMSa}{"02}

%
% Macros for moduli stacks/spaces
%
\def\QCohstack{\mathcal{QC}\!\mathit{oh}}
\def\Cohstack{\mathcal{C}\!\mathit{oh}}
\def\Spacesstack{\mathcal{S}\!\mathit{paces}}
\def\Quotfunctor{\mathrm{Quot}}
\def\Hilbfunctor{\mathrm{Hilb}}
\def\Curvesstack{\mathcal{C}\!\mathit{urves}}
\def\Polarizedstack{\mathcal{P}\!\mathit{olarized}}
\def\Complexesstack{\mathcal{C}\!\mathit{omplexes}}
% \Pic is the operator that assigns to X its picard group, usage \Pic(X)
% \Picardstack_{X/B} denotes the Picard stack of X over B
% \Picardfunctor_{X/B} denotes the Picard functor of X over B
\def\Pic{\mathop{\mathrm{Pic}}\nolimits}
\def\Picardstack{\mathcal{P}\!\mathit{ic}}
\def\Picardfunctor{\mathrm{Pic}}
\def\Deformationcategory{\mathcal{D}\!\mathit{ef}}
\usepackage{cleveref}
\crefname{originalreference}{original reference}{original references}
\Crefname{originalreference}{Original reference}{Original references}
\makeatletter
\newlabel{situation-quasi-compact-objects}{{0937}{source}{}{sourcepointer-1}{}}
\newlabel{situation-quasi-compact-objects@cref}{{[originalreference][0][]0937}{[1][0][]source}{}{}{}}
\makeatother
\crefname{originalreference}{original reference}{original references}
\Crefname{originalreference}{Original reference}{Original references}
\makeatletter
\newlabel{item-enough}{{0938}{source}{}{sourcepointer-2}{}}
\newlabel{item-enough@cref}{{[originalreference][0][]0938}{[1][0][]source}{}{}{}}
\makeatother
\crefname{originalreference}{original reference}{original references}
\Crefname{originalreference}{Original reference}{Original references}
\makeatletter
\newlabel{item-enough-qc}{{0939}{source}{}{sourcepointer-3}{}}
\newlabel{item-enough-qc@cref}{{[originalreference][0][]0939}{[1][0][]source}{}{}{}}
\makeatother
\crefname{originallemma}{original lemma}{original lemmas}
\Crefname{originallemma}{Original lemma}{Original lemmas}
\makeatletter
\newlabel{lemma-module-quotient-direct-sum}{{093B}{source}{}{sourcepointer-4}{}}
\newlabel{lemma-module-quotient-direct-sum@cref}{{[originallemma][0][]093B}{[1][0][]source}{}{}{}}
\makeatother
\crefname{originalreference}{original reference}{original references}
\Crefname{originalreference}{Original reference}{Original references}
\makeatletter
\newlabel{sites-lemma-directed-colimits-sections}{{0738}{source}{}{sourcepointer-5}{}}
\newlabel{sites-lemma-directed-colimits-sections@cref}{{[originalreference][0][]0738}{[1][0][]source}{}{}{}}
\makeatother
\crefname{originallemma}{original lemma}{original lemmas}
\Crefname{originallemma}{Original lemma}{Original lemmas}
\makeatletter
\newlabel{lemma-quasi-compact-hom-from}{{0936}{source}{}{sourcepointer-6}{}}
\newlabel{lemma-quasi-compact-hom-from@cref}{{[originallemma][0][]0936}{[1][0][]source}{}{}{}}
\makeatother
\title{Sheaves and modules as filtered colimits of standard presentations built from a basis of quasi-compact objects}
\author{The Stacks Project Authors}
\date{stacks-project:sites-modules.tex:local-snapshot-20261012; GNU FDL 1.2 or later}
\begin{document}
\maketitle
\noindent\textit{Editorial note.} All delivered mathematical text is the author's own, copied verbatim from the pinned Stacks Project chapter file sites-modules.tex; no statement or proof step was rewritten, reordered or removed.\par\medskip
\noindent\textit{Editorial note.} Both the statement and the proof contain commutative diagrams typeset with Xy-pic; the retained two-slice preamble of this source supplies the native Xy-pic setup, so they compile as in the source.\par\medskip
\noindent\textit{Editorial note.} The counted span cites a situation and a quotient lemma, both recorded in assembly.external\_references with canonical tags as navigation pointers.\par\medskip
\noindent\textit{Allowed standard prerequisites (editorial).}\begin{itemize}
\item Every sheaf of sets on a site with enough quasi-compact objects is a quotient of a coproduct of representables. Stacks Project, Modules on Sites, Lemma lemma-module-quotient-direct-sum (tag 09B3). Produces the surjection and the fibre product presentation used to build the coequalizer.
\item The quasi-compact objects situation. Stacks Project, Modules on Sites, Situation situation-quasi-compact-objects (tag 09AZ). Supplies the basis B and the quasi-compactness hypotheses that make the finite cutting argument work.
\end{itemize}\medskip
\ifdefined\hypertarget\hypertarget{sourcepointer-1}{}\fi
\noindent\textit{Source reference.} Original reference 0937: sources/stacks-chow/tags.txt, source lines 11505--11505. Cross-chapter standard prerequisite cited by the author inside the counted span; it is used as a pointer and its own proof is not part of this item.\par\medskip
\ifdefined\hypertarget\hypertarget{sourcepointer-2}{}\fi
\noindent\textit{Source reference.} Original reference 0938: sources/stacks-chow/tags.txt, source lines 11506--11506. Cross-chapter standard prerequisite cited by the author inside the counted span; it is used as a pointer and its own proof is not part of this item.\par\medskip
\ifdefined\hypertarget\hypertarget{sourcepointer-3}{}\fi
\noindent\textit{Source reference.} Original reference 0939: sources/stacks-chow/tags.txt, source lines 11507--11507. Cross-chapter standard prerequisite cited by the author inside the counted span; it is used as a pointer and its own proof is not part of this item.\par\medskip
\ifdefined\hypertarget\hypertarget{sourcepointer-4}{}\fi
\noindent\textit{Source reference.} Original lemma 093B: sources/stacks-chow/tags.txt, source lines 11509--11509. Cross-chapter standard prerequisite cited by the author inside the counted span; it is used as a pointer and its own proof is not part of this item.\par\medskip
\ifdefined\hypertarget\hypertarget{sourcepointer-5}{}\fi
\noindent\textit{Source reference.} Original reference 0738: sources/stacks-chow/tags.txt, source lines 9052--9052. Cross-chapter standard prerequisite cited by the author inside the counted span; it is used as a pointer and its own proof is not part of this item.\par\medskip
\ifdefined\hypertarget\hypertarget{sourcepointer-6}{}\fi
\noindent\textit{Source reference.} Original lemma 0936: sources/stacks-chow/tags.txt, source lines 11504--11504. Cross-chapter standard prerequisite cited by the author inside the counted span; it is used as a pointer and its own proof is not part of this item.\par\medskip
\begin{lemma}
\label{lemma-module-filtered-colimit-constructibles}
In Situation \ref{situation-quasi-compact-objects} assume
(\ref{item-enough}) and (\ref{item-enough-qc}) hold.
\begin{enumerate}
\item Every sheaf of sets is a filtered colimit of sheaves of the form
\begin{equation}
\label{equation-towards-constructible-sets}
\text{Coequalizer}\left(
\xymatrix{
\coprod\nolimits_{j = 1, \ldots, m} h_{V_j}^\#
\ar@<1ex>[r] \ar@<-1ex>[r] &
\coprod\nolimits_{i = 1, \ldots, n} h_{U_i}^\#
}
\right)
\end{equation}
with $U_i$ and $V_j$ in $\mathcal{B}$.
\item If $\mathcal{O}$ is a sheaf of rings, then every $\mathcal{O}$-module
is a filtered colimit of sheaves of the form
\begin{equation}
\label{equation-towards-constructible}
\Coker\left(
\bigoplus\nolimits_{j = 1, \ldots, m} j_{V_j!}\mathcal{O}_{V_j}
\longrightarrow
\bigoplus\nolimits_{i = 1, \ldots, n} j_{U_i!}\mathcal{O}_{U_i}
\right)
\end{equation}
with $U_i$ and $V_j$ in $\mathcal{B}$.
\end{enumerate}
\end{lemma}

\begin{proof}
Proof of (1). By Lemma \ref{lemma-module-quotient-direct-sum}
every sheaf of sets $\mathcal{F}$ is the target of a surjection
whose source is a coprod $\mathcal{F}_0$
of sheaves the form $h_{U}^\#$ with $U \in \mathcal{B}$.
Applying this to $\mathcal{F}_0 \times_\mathcal{F} \mathcal{F}_0$
we find that $\mathcal{F}$ is a coequalizer of a pair of maps
$$
\xymatrix{
\coprod\nolimits_{j \in J} h_{V_j}^\#
\ar@<1ex>[r] \ar@<-1ex>[r] &
\coprod\nolimits_{i \in I} h_{U_i}^\#
}
$$
for some index sets $I$, $J$ and $V_j$ and $U_i$ in $\mathcal{B}$.
For every finite subset $J' \subset J$ there is a finite subset
$I' \subset I$ such that the coproduct over $j \in J'$ maps into
the coprod over $i \in I'$ via both maps, see
Sites, Lemma \ref{sites-lemma-directed-colimits-sections}.
(Details omitted; hint: an infinite coproduct is the filtered colimit
of the finite sub-coproducts.)
Thus our sheaf is the colimit of the cokernels of these maps
between finite coproducts.

\medskip\noindent
Proof of (2).
By Lemma \ref{lemma-module-quotient-direct-sum}
every module is a quotient of a direct sum of modules of the form
$j_{U!}\mathcal{O}_U$ with $U \in \mathcal{B}$. Thus every module
is a cokernel
$$
\Coker\left(
\bigoplus\nolimits_{j \in J} j_{V_j!}\mathcal{O}_{V_j}
\longrightarrow
\bigoplus\nolimits_{i \in I} j_{U_i!}\mathcal{O}_{U_i}
\right)
$$
for some index sets $I$, $J$ and $V_j$ and $U_i$ in $\mathcal{B}$.
For every finite subset $J' \subset J$ there is a finite subset
$I' \subset I$ such that the direct sum over $j \in J'$ maps into
the direct sum over $i \in I'$, see
Lemma \ref{lemma-quasi-compact-hom-from}.
Thus our module is the colimit of the cokernels of these maps
between finite direct sums.
\end{proof}

\end{document}
